
\section{在辛几何中的应用}
\label{sec:1-3}

本节说明弗洛尔上同调在拉格朗日子流形的辛拓扑中的应用。我们先陈述由定理 B、C、D 和 G 合并得到的如下推论。

\begin{theorem}{H}
设 $L$ 是相对 spin 的拉格朗日子流形，并假定自然映射
$H_*(L;\mathbb Q)\to H_*(M;\mathbb Q)$ 是单射。则对任意使
$L$ 与 $\phi(L)$ 横截的哈密顿微分同胚 $\phi:M\to M$，有
\[
  \#(L\cap\phi(L))\geq \sum_k\operatorname{rank}H_k(L;\mathbb Q).
\]
\end{theorem}

定理 H 的证明见第 6.5.1 小节。注意，对角子流形 $\Delta\subset M\times M$ 满足定理 H 的假设，因而 [FuOn99II, LiuTi98, Rua99] 针对哈密顿微分同胚（取 $\mathbb Q$ 系数）证明的阿诺尔德猜想\liftIndex{I001}版本可由定理 H 推出。下面给出定理 H 的一个更精确版本。

\begin{theorem}{I}
设 $L$ 相对 spin，并假定 $\mathcal M_{\mathrm{weak,def}}(L)$ 非空。记
\[
  A=\sum_*\operatorname{rank}H_*(L;\mathbb Q),\qquad
  B=\sum_*\operatorname{rank}\operatorname{Ker}
  \bigl(H_*(L;\mathbb Q)\to H_*(M;\mathbb Q)\bigr).
\]
则对任意使 $L$ 与 $\phi(L)$ 横截的哈密顿微分同胚
$\phi:M\to M$，有
\[
  \#(L\cap\phi(L))\geq A-2B.
\]
\end{theorem}

定理 I 的证明在第 6.5.2 小节完成。如前所述，取
$\Lambda_{0,\mathrm{nov}}$ 系数的弗洛尔上同调会随拉格朗日子流形对
$L_1,L_0$ 的哈密顿同痕而改变，而取 $\Lambda_{\mathrm{nov}}$ 系数时则不随之改变。我们更仔细地研究这一现象，并把弗洛尔上同调的“挠部分”\liftIndex{I016}\liftIndex{I033}同哈密顿微分同胚之间的霍费尔距离联系起来。（霍费尔距离的定义见 [HoZe94]。）首先注意，弗洛尔上同调一般具有分解
\begin{equation}
  HF\bigl((L_1,b_1),(L_0,b_0);\Lambda_{0,\mathrm{nov}}\bigr)
  \cong
  \bigl(\Lambda_{0,\mathrm{nov}}\bigr)^a
  \oplus
  \bigoplus_{i=1}^{b}
  \frac{\Lambda_{0,\mathrm{nov}}}
       {T^{\lambda_i}\Lambda_{0,\mathrm{nov}}}.
  \tag{1.13}\label{eq:1-13}
\end{equation}
这里用到了定理 6.1.20。我们称 $a\in\mathbb Z_{\geq0}$ 为贝蒂数，称
$\lambda_i\in(0,\infty)$ 为弗洛尔上同调的挠指数。（如果 $L_1$ 与
$L_0$ 干净相交，则整数 $b$ 有限；一般情况下它也可能无穷，见例 6.5.40。）

\begin{theorem}{J}
设 $(L_1,L_0)$ 是 $M$ 中相对 spin 的一对拉格朗日子流形，并且
\[
  (b_1,b_0)\in
  \mathcal M_{\mathrm{weak,def}}(L_1)
  \mathop{\times}_{(\pi_{\mathrm{amb}},\mathfrak{PO})}
  \mathcal M_{\mathrm{weak,def}}(L_0).
\]
设 $\phi:M\to M$ 是哈密顿微分同胚，$\mu$ 是它的霍费尔范数，即
$\mu=\lVert\phi\rVert$。假定 $\phi(L_1)$ 与 $L_0$ 横截，并记
\[
  b(\mu)=\#\{i\mid\lambda_i\geq\mu\},
\]
其中 $\lambda_i$ 是式 \eqref{eq:1-13} 中的挠指数。则
\[
  \#\bigl(\phi(L_1)\cap L_0\bigr)\geq a+2b(\mu).
\]
\end{theorem}

定理 J 的证明在第 6.5.3 小节完成。定理 J 与切卡诺夫的工作 [Chek96,98] 密切相关；我们在证明中采用了切卡诺夫的想法。事实上，定理 J 的一个小幅加强（第 6.5.5 小节的定理 6.5.47）蕴含切卡诺夫定理。在第 6.5.3 小节中，我们证明了比定理 J 更强的结果：当 $L_i$ 在其哈密顿同痕类中移动时，各挠指数 $\lambda_i$ 都是关于霍费尔距离的 $1$-利普希茨函数（定理 6.1.25）。

下一个应用涉及拉格朗日嵌入 $L\subset\mathbb C^n$ 的马斯洛夫类\liftIndex{I022}
$\mu_L\in H^1(L;\mathbb Z)$。一个广为流传的猜想断言，任意紧拉格朗日嵌入到
$\mathbb C^n$ 中所得的马斯洛夫类都非平凡。关于这一猜想已有若干部分结果（见
[Vit90, Pol91I, Oh96I]）；这里我们再给出一个。

\begin{theorem}{K}
设 $L$ 是 $\mathbb C^n$ 中满足 $H^2(L;\mathbb Q)=0$ 的紧 spin 拉格朗日子流形，则它的马斯洛夫类 $\mu_L\in H^1(L;\mathbb Z)$ 非零。
\end{theorem}

定理 K 由第 6.4.4 小节证明的推论 6.1.16 得到。定理 K 可以推广到更一般的
$(M,L)$；参见定理 L 和推论 6.1.16。利用谱序列同乘法结构的相容性（定理 D 的
(D.4)），还可得到另一个应用，即定理 6.4.35。

\begin{theorem}{L}
设 $L\subset(M,\omega)$ 是相对 spin 的拉格朗日子流形，并且马斯洛夫指标同态
\[
  I_\mu:\pi_2(M,L)\longrightarrow\mathbb Z
\]
平凡。再假定 $H^2(L;\mathbb Q)=0$。则对任意哈密顿微分同胚
$\phi:M\to M$，均有 $L\cap\phi(L)\neq\varnothing$。此外，若 $L$ 与
$\phi(L)$ 横截，则存在马斯洛夫指标为 $0$ 的交点
$p\in L\cap\phi(L)$。
\end{theorem}

定理 L 在第 6.4.3 小节证明。满足马斯洛夫指标同态
$I_\mu:\pi_2(M,L)\to\mathbb Z$ 平凡的一类特殊对子 $(M,L)$ 如下：
$(M,\omega,J)$ 是卡拉比--丘流形\liftIndex{I011}，而 $L\subset M$ 是马斯洛夫类平凡的拉格朗日子流形。在这种情形下，马斯洛夫指标 $\mu_L(\beta)$ 只依赖于
$\partial\beta\in\pi_1(L)$，因而诱导出 $H^1(L;\mathbb Z)$ 中一个良定义的上同调类，称为 $L$ 的马斯洛夫类。（进一步说明见 [Mor81, Daz81, HaLa82]。）理查德·托马斯（Richard P. Thomas）与丘成桐（Shing-Tung Yau）随后在 [ThYa02] 中用定理 L 证明了一个唯一性定理：在卡拉比--丘三维流形中，特殊拉格朗日同调三维球面（即满足 $H^1(L;\mathbb Z)=0$ 的 $L$）在其哈密顿同痕类中是唯一的。

如前所述，弗洛尔上同调
$HF((L,b),(L,b);\Lambda_{0,\mathrm{nov}})$ 不一定同构于 $L$ 的上同调；两者之差由定理 D 中谱序列的微分 $\delta_r$ 决定，而这些微分又由适当的伪全纯圆盘模空间的基本链决定。因此，要实际计算
$HF((L,b),(L,b);\Lambda_{0,\mathrm{nov}})$，一般很难分析这个谱序列，因为这需要仔细研究伪全纯圆盘的计数问题\liftIndexRangeStart{I045}。这个计数问题甚至比通常格罗莫夫--威腾理论中的伪全纯球面计数更为精细而困难。另一方面，定理 A 表明，拉格朗日子流形的通常上同调和弗洛尔上同调上存在各种可用于计算的乘积结构；但在具体例子中计算这些乘积一般也并非易事。下面列出若干计算。

\medskip
\noindent\textup{(1.14.0)}\quad
若 $\pi_2(M,L)=\pi_1(L)=0$，则定理 X 表明，定理 A 中的滤过
$A_\infty$ 代数无非就是 $L$ 的实同伦型。

\medskip
\noindent\textup{(1.14.1)}\quad
Cheol-Hyun Cho 与第二作者吴龙根（Yong-Geun Oh）针对法诺环面流形\liftIndex{I031}的拉格朗日环面轨道完成了一些计算 [Cho04I, ChOh03]。他们的工作部分受到堀健太朗（Kentaro Hori；日文 堀健太朗）与瓦法的工作 [HoVa00, Hor01] 启发。特别地，[ChOh03] 计算了定理 B 中函数
\[
  \mathfrak{PO}:\mathcal M_{\mathrm{weak,def}}(L)
  \longrightarrow\Lambda_{0,\mathrm{nov}}^+
\]
的一个变体；它对应于物理学中的朗道--金兹堡超势 [HoVa00]。粗略地说，他们证明了：弗洛尔上同调非平凡的环面轨道通过镜像映射对应于该超势的临界点。在 [ChOh03] 中，Cheol-Hyun Cho 与吴龙根称这些轨道为平衡环面纤维\liftIndex{I032}。

[ChOh03] 还证明，平衡环面纤维的弗洛尔上同调 $HF(L,L)$ 同构于经典上同调群
$H^*(T^n;\Lambda_{\mathrm{nov}})$。在 [Cho04II] 中，Cheol-Hyun Cho 证明由
$A_\infty$ 结构在 $HF(L,L)$ 上诱导的自然环结构，经朗道--金兹堡镜像映射后，同构于该超势的黑塞矩阵所确定的二次型相伴的克利福德代数。堀健太朗 [Hor01] 以及安东·卡普斯汀（Anton Kapustin）与 Yi Li [KapLi04] 此前已经预言了这一点。特别地，这一计算给出了一个例子：弗洛尔上同调同构于经典上同调，但其环结构是经典杯积的变形。回忆环面的上同调环同构于由其一次上同调生成的外代数。

\medskip
\noindent\textup{(1.14.2)}\quad
[FOOO08I, FOOO08II] 使用并推广了 [ChOh03] 的结果；本书作者在那里应用本书的许多结果，研究紧环面流形（不必是法诺流形）中以 $T^n$ 为轨道的拉格朗日子流形的弗洛尔上同调。

\medskip
\noindent\textup{(1.14.3)}\quad
若 $L\subseteq\mathbb C^{n+1}$ 是同伦等价于 $S^1\times S^n$ 的拉格朗日子流形，则 [FOOO09I] 详细研究了相应的谱序列。

\medskip
\noindent\textup{(1.14.4)}\quad
[FOOO09II] 对如下情形作了一些计算：$L$ 是由 $T^{2n}$ 中的平坦拉格朗日子环面经拉格朗日手术得到的。（另见 [Fuk02III]。）Theta 函数出现在这一计算中；[FOOO09II] 还给出了其他许多例子。

\medskip
\noindent\textup{(1.14.5)}\quad
同调镜像对称猜想\liftIndex{I019}一旦得到证明，就会成为计算拉格朗日子流形的弗洛尔上同调和滤过 $A_\infty$ 代数的有力工具。例如，保罗·塞德尔（Paul Seidel）关于四次曲面同调镜像对称的证明 [Sei03II]，应当会对这些结构的计算给出一些推论。

\medskip
\noindent\textup{(1.14.6)}\quad
物理学文献中有大量结果研究附着在给定拉格朗日子流形上的全纯圆盘计数问题\liftIndexRangeEnd{I045}，尤其是在环境辛流形为非紧卡拉比--丘流形时，这些结果与局部镜像原理有关（见 [AAMV05, DiFl05, LLLZ04] 及其中所引文献）。不过，目前还不清楚如何把这些研究用于实际计算本书构造的结构。
